% 本文件由 LaTeX 排版 Agent 根据有效 translation.json 自动生成。
% author=Origen; work_id=0415; title=致额我略书

\chapter{\WorkTitle{致额我略书}}

% structure: p0001 source=h1 target=omitted reason=page-title-already-rendered-by-parent

1. 来自奥利振的问候：最卓越的先生，可敬的儿子额我略，愿你安好于天主。天资聪颖，如你所知，若加以操练，或许能对那所谓偶然的目的有所贡献——即任何人心所愿之事。因此，你的良好天赋本可使你成为一位完美的罗马法学家或希腊哲学家，姑且说，成为备受推崇的某个学派的代表。但我切望你将全部天赋之力奉献于基督信仰，以之为你的目的；为此，我愿你从希腊人的哲学中提取那可作为基督信仰的课程或预备之物，并从几何学和天文学中提取那有助于解释圣经的学问，如此，哲学的儿子们惯常论及几何学、音乐、文法、修辞学与天文学作为哲学的辅助者，我们也可对哲学本身这样说，相对于基督信仰。\par

2. 或许这类事在\BibleBook{出谷}{纪}中从天主口中所写的话里已有预表：以色列子民奉命向他们的邻人及与他们同住的人索取金银器皿和衣裳，为的是藉掳掠埃及人，他们得以有材料制备与敬礼天主有关之物。因为从以色列子民取自埃及人的东西中，制成了至圣所中的器皿——约柜及柜盖、革鲁宾、赎罪盖、以及盛有玛纳——天使之粮——的金盒。这些大概是用最好的埃及金子制成的。次等的金子则用于靠近内帐幔的纯金灯台及其枝子、放陈设饼的金桌、以及其间的金香炉。若有第三、第四等的金子，则用以制造圣器；其余之物用埃及的银子制作。因为以色列子民居住在埃及时，从那里的居住获得了这等好处：他们不缺如此宝贵的材料来制作敬礼天主的器皿。而那一切需要缝制和刺绣的工作，据圣经记载，是藉天主的智慧彼此缝制，以制成帐幔、内院和外院，大概也是用埃及的衣裳做的。我何必在这不合时宜的离题话中继续说明，从埃及取来的东西对以色列子民多么有益——那些东西埃及人没有善用，而希伯来人在天主的智慧引导下却用以事奉天主？圣经通常将以色列子民从本地下降到埃及描述为一种恶，意指某些人在与埃及人交往（即涉足世俗学问）之后，受到损害，尽管他们已遵从了天主的法律和以色列的事奉。至少，厄东人阿德尔就是一例：他留在以色列地、未尝埃及之饼时，并未制造偶像。但当他逃离智慧的撒罗满，下到埃及，仿佛逃离天主的智慧，并因娶了法郎妻子的妹妹而成为法郎的亲属，生养儿女与法郎的子女一同长大时，他就做了这事。因此，他虽然回到以色列地，却只为了分裂天主的子民，并使他们向金牛犊说：\BibleQuote{这些是你的神，以色列，曾领你出埃及地。}我可以凭我的经验告诉你，没有多少人只从埃及取用有益之物，然后离开并用以事奉天主；而厄东人阿德尔却有许多弟兄。这些人从他们的希腊研究中产生异端观念，并将它们如同金牛犊一样设立在贝特耳——这名字意为\Quote{天主的殿}。在我看来，这些话也表明：他们在圣经中——那里是天主圣言居住之处，象征性地称为贝特耳——设立了己意。另一个像，经文说，设立在丹。丹的边界是最边远的，靠近外邦人的边界，正如在\BibleBook{若苏厄}{书}中所写的那样。这些阿德尔的弟兄们（我们如此称呼）的某些发明也很靠近外邦人的边界。\par

3. 因此，我儿，你要勤奋地阅读圣经。我说，你要勤奋。因为我们这些阅读天主之事的人，需要很大的勤奋，以免对它们说或想得过于草率。当你这样勤奋地研究天主之事时，要怀着忠信的预断（即那些中悦天主的预断），叩那锁着的门，它必为你开启，由那门房——耶稣说：\BibleQuote{给他开门的就是那门房。}当你这样研读圣经时，要正直地寻求，并坚定地信赖天主，寻求那许多人所错失的圣经意义。不要满足于叩门和寻求；因为祈祷对于认识天主之事是必不可少的。为此救主劝勉说，不仅说：\BibleQuote{叩，就给你们开门；寻求，就必寻见，}而且说：\BibleQuote{求，就必给你们。}我作为父亲对你的爱使我如此大胆；但我的大胆是否良好，天主和祂的基督，以及所有分沾天主圣神和基督圣神的人，都将知道。愿你也能分沾，并不断增加你的基业，以便你不仅说：\Quote{我们成了基督的分享者，}而且成了天主的分享者。\par

\SourceRecord{Translated by Frederick Crombie.；Ante-Nicene Fathers；Vol. 4.；Edited by Alexander Roberts, James Donaldson, and A. Cleveland Coxe.；Buffalo, NY: Christian Literature Publishing Co.,；1885.；<http://www.newadvent.org/fathers/0415.htm>.}
